\section{Mocking}
\linksubsection{unittest-mock}{unittest.mock}

The \mono{pytest-mock} plugin exposes \mono{unittest.mock} features via a \mono{mocker} fixture.
\mono{Mock} objects are \enquote{magic} objects that can pretend to be anything.
\vspace{-.5em}

\begin{minted}[autogobble]{python}
    >>> dev = Mock(name="dev")
\end{minted}
\vspace{-.5em}

On a \mono{Mock} object, we can access any attributes, call the mock, etc.
We always get another \mono{Mock} back:
\vspace{-.25em}

\begin{minted}[autogobble]{python}
    >>> dev.conn
    <Mock name='dev.conn' id='X'>
\end{minted}
\vspace{-1.5em}

\begin{minted}[autogobble]{python}
    >>> dev.conn.send("cmd")
    <Mock name='dev.conn.send()' id='...4736'>
\end{minted}
\vspace{-.25em}

\newenvironment{mocktable}{
    \begin{tabular}{>{\color{bruhinblue}}l>{\ttfamily}l}
        }{%
    \end{tabular}
}

\vspace{1em}
\begin{center}
    \begin{tikzpicture}[
            mock/.style={
                    bluenodebright,
                    align=left,
                    label={[label distance=.5em]left:\mono{Mock}},
                    minimum width=5.5cm,
                    text width=5cm,
                },
        ]
        \node [mock] (dev) {
            \begin{mocktable}
                name & dev             \\
                id   & \codeskip{}1376 \\
            \end{mocktable}
        };

        \node [mock, below=8mm of dev, bluenode] (conn) {
            \begin{mocktable}
                name & dev.conn \\
                id   & \emph{X} \\
            \end{mocktable}
        };
        \draw [->] (dev) -- (conn) node[midway, right] {\mono{.conn}};

        \node [mock, below=8mm of conn] (send) {
            \begin{mocktable}
                name & dev.conn.send   \\
                id   & \codeskip{}3056 \\
            \end{mocktable}
        };
        \draw [->] (conn) -- (send) node[midway, right] {\monot{\codeskip.send}};

        \node [mock, below=8mm of send] (call) {
            \begin{mocktable}
                name & dev.conn.send() \\
                id   & \codeskip{}4736 \\
            \end{mocktable}
        };
        \draw [->] (send) -- (call) node[midway, right] {\monot{\codeskip()}};
    \end{tikzpicture}
\end{center}
\vspace{1em}

Doing the same thing again results in the same mock (same \mono{id}):
\vspace{-.25em}

\begin{minted}[autogobble]{python}
    >>> dev.conn
    <Mock name='dev.conn' id='X'>
\end{minted}
\pagebreak

We can set special attributes to \emph{configure} the mock's behavior:

\begin{minted}[autogobble,escapeinside=||,beameroverlays]{python}
    >>> dev.conn.|\highlightbox[white]{send}|.return_value = "ok"
    >>> dev.conn.|\highlightbox{send(\PYG{s}{\PYGZdq{}cmd\PYGZdq{}})}|
    'ok'
\end{minted}

Then we run the tests, which use the mock like a real object. The mock remembers what happened, so we can assert on that:

\begin{minted}[autogobble,escapeinside=||,beameroverlays]{python}
    >>> dev.conn.|\highlightbox{send.assert\_called\_with(\PYG{s}{\PYGZdq{}cmd\PYGZdq{}})}|
\end{minted}

\vspace{-.5em}
\begin{minted}[autogobble,escapeinside=||]{python}
    >>> dev.conn.|\highlightbox[highlightcolorred]{send.assert\_called\_with(\PYG{s}{\PYGZdq{}no\PYGZdq{}})}|
    |\codeskip|
    AssertionError: |expected call not found.|
    Expected: |\highlightbox[highlightcolorred]{send(\PYG{s}{\PYGZsq{}no\PYGZsq{}})}|
    Actual:   |\highlightbox{send(\PYG{s}{\PYGZsq{}cmd\PYGZsq{}})}|
\end{minted}

\subsection{Mocking libraries}

Sometimes there are ready-made mocking libraries, so there is no need to hand-roll mocks:

\begin{description}[labelwidth=9em, leftmargin=!, itemsep=0pt]
    \item[requests] \textbf{responses}, betamax, \\ requests-mock
    \item[aiohttp] aioresponses
    \item[HTTP in general] \textbf{VCR.py}, HTTPretty, pook
    \item[sockets in general] mocket
    \item[time/datetime] freezegun, time-machine
\end{description}

For most of those, there's a corresponding pytest plugin for nice integration.

\textbf{Result:} More high-level abstraction for mocking \\
(e.g.\ HTTP responses instead of the Python API).

Well-tested mocks, usually covering the entire API of a library.

Not many high-level mocking libraries exist. \\ But if one does, \textbf{use it}!
